% ============ 第4章 §4.6 Fourier限制问题(Tomas-Stein定理)(原书 pp.338-346) ============
\section{Fourier 限制问题(Tomas--Stein 定理)}\label{sec:4.6}
\index{Fourier 限制问题}
\index{Tomas--Stein 定理}


在上一节中,利用在 $L^{2}$ 估计中可以直接写出其范数这一点,对振荡积分算子进行了估计.本节将进一步深入讨论,考察基于 $L^{2}$ 理论的 Fourier 限制问题.

由系~\ref{cor:4.3.10},$f \in L^{p}(\mathbb{R}^{n})$($1 \leq p \leq 2$)的 Fourier 变换 $\hat{f}$ 是函数,且满足
\begin{equation}
\|\hat{f}\|_{p'} \leq \|f\|_{p} \quad \left(p' = \frac{p}{p-1}\right).\tag{4.6.1}
\label{eq:4.6.1}\end{equation}

$L^{1}(\mathbb{R}^{n})$ 的 Fourier 像是既有界又连续的函数,故将 Fourier 像整体限制在 $\mathbb{R}^{n}$ 的任意子集上都是有意义的.另一方面,$L^{2}(\mathbb{R}^{n})$ 的 Fourier 像是整个 $L^{2}(\mathbb{R}^{n})$,故将 Fourier 像整体限制在 Lebesgue 测度为 0 的集合上没有意义.那么,处于这两者之间的 $L^{p}(\mathbb{R}^{n})$($1 < p < 2$)是如何反映这一极端差异的呢?

令 $\phi \in \mathcal{S}(\mathbb{R}^{n-1})$,
\begin{equation*}
f(x_{1}, x_{2}, \ldots, x_{n}) := \frac{1}{1 + |x_{1}|}\phi(x_{2}, \ldots, x_{n}),
\end{equation*}
则 $f \in L^{p}(\mathbb{R}^{n})$($p > 1$),且在超平面 $\{\xi_{1} = 0\}$ 上 $\hat{f}(\xi) = \infty$.但是,1967 年 E. M. Stein 证明了:在某个 $p$ 的范围内,可以将 $L^{p}(\mathbb{R}^{n})$ 的 Fourier 像限制到具有某种曲率的超曲面上.当然,由 \eqref{eq:4.6.1},限制到正测度集合上总是可能的.

这里讨论单位球面
\begin{equation*}
S^{n-1} = \{x \in \mathbb{R}^{n}: |x| = 1\}
\end{equation*}
上的 Fourier 限制问题.

证明下面的定理是本节的目标.以下用 $d\sigma$ 表示定义在 $S$ 上的 Lebesgue 测度.

\begin{theorem}[Tomas--Stein]\label{thm:4.6.1}
\begin{equation}
\left(\int_{S}|\hat{f}(\xi)|^{q}\dd\sigma\right)^{1/q} \leq C_{p}\|f\|_{L^{p}(\mathbb{R}^{n})} \quad (f \in \mathcal{S}(\mathbb{R}^{n}))\tag{4.6.2}
\label{eq:4.6.2}\end{equation}
当
\begin{equation*}
q = \frac{n-1}{n+1}\frac{p}{p-1}
\end{equation*}
时,在
(1) $1 \leq p \leq \frac{2(n+1)}{n+3}$,$n \geq 3$,
(2) $1 \leq p < \frac{4}{3}$,$n = 2$
的情形下成立.
\end{theorem}

令
\begin{equation*}
\begin{cases}
D := \{y \in \mathbb{R}^{n-1}: |y| \leq 1/2\}, \\
h(y) := \sqrt{1 - |y|^{2}} \quad (y \in D), \\
\tilde{D} := \{(y, h(y)) \in S: y \in D\},
\end{cases}
\end{equation*}
则在 $\tilde{D}$ 上
\begin{equation*}
\dd\sigma = \sqrt{1 + |\nabla h(y)|^{2}}\dd y = \frac{\dd y}{h(y)}.
\end{equation*}
考虑对称性与光滑的单位分解,为证明 \eqref{eq:4.6.2},固定 $0 \leq b(y) \in C_{0}^{\infty}(D)$,只要证明
\begin{equation*}
\left(\int_{\mathbb{R}^{n-1}}|\hat{f}(y, h(y))|^{q}a(y)\dd y\right)^{1/q} \leq C\|f\|_{L^{p}(\mathbb{R}^{n})}, \quad a(y) := \frac{b(y)}{h(y)}
\end{equation*}
即可.再考虑对偶性,证明
\begin{equation}
\left|\int_{\mathbb{R}^{n-1}}\hat{f}(y, h(y))g(y)a(y)\dd y\right| \leq C\|f\|_{L^{p}(\mathbb{R}^{n})}\times\|g\|_{L^{q'}(\mathbb{R}^{n-1})}.\tag{4.6.3}
\label{eq:4.6.3}\end{equation}

由于
\begin{equation}
\begin{aligned}
& \int_{\mathbb{R}^{n-1}}\hat{f}(y, h(y))g(y)a(y)\dd y \\
& = \int_{\mathbb{R}^{n-1}}\int_{\mathbb{R}^{n}}e^{-2\pi i(z_{1}y_{1} + \ldots + z_{n-1}y_{n-1} + z_{n}h(y))}f(z)g(y)a(y)\dd z\dd y,
\end{aligned}
\end{equation}
\eqref{eq:4.6.3} 归结为证明
\begin{equation}
\left\|\int_{\mathbb{R}^{n-1}}e^{-2\pi i(z_{1}y_{1} + \ldots + z_{n-1}y_{n-1} + z_{n}h(y))}g(y)a(y)\dd y\right\|_{L^{p'}(\mathbb{R}^{n})} \leq C\|g\|_{L^{q'}(\mathbb{R}^{n-1})}.\tag{4.6.4}
\label{eq:4.6.4}\end{equation}

用上一节的振荡积分算子方法分析 \eqref{eq:4.6.4}.即设 $a(z, y) \in C_{0}^{\infty}(\mathbb{R}^{n}\times\mathbb{R}^{n-1})$,证明
\begin{equation}
\begin{aligned}
& \left\|\int_{\mathbb{R}^{n-1}}e^{-2\pi i\lambda(z_{1}y_{1} + \ldots + z_{n-1}y_{n-1} + z_{n}h(y))}a(z, y)g(y)\dd y\right\|_{L^{p'}(\mathbb{R}^{n})} \\
& \leq C\lambda^{-n/p'}\|g\|_{L^{q'}(\mathbb{R}^{n-1})}
\end{aligned}\tag{4.6.5}
\label{eq:4.6.5}\end{equation}
成立.只要有这个式子,变形为
\begin{equation}
\begin{aligned}
& \left\|\int_{\mathbb{R}^{n-1}}e^{-2\pi i(z_{1}y_{1} + \ldots + z_{n-1}y_{n-1} + z_{n}h(y))}a(z/\lambda, y)g(y)\dd y\right\|_{L^{p'}(\mathbb{R}^{n})} \\
& \leq C\|g\|_{L^{q'}(\mathbb{R}^{n-1})},
\end{aligned}
\end{equation}
令 $\lambda \to \infty$,即得 \eqref{eq:4.6.4}.以下,令
\begin{equation*}
\phi(z, y) := z_{1}y_{1} + \ldots + z_{n-1}y_{n-1} + z_{n}h(y),
\end{equation*}
考察振荡积分算子
\begin{equation*}
T_{\lambda}g(z) := \int_{\mathbb{R}^{n-1}}e^{i\lambda\phi(z, y)}a(z, y)g(y)\dd y.
\end{equation*}

整理要点,证明下面的命题.

\begin{proposition}\label{prop:4.6.2}
\begin{equation}
\|T_{\lambda}g\|_{L^{p'}(\mathbb{R}^{n})} \leq C\lambda^{-n/p'}\|g\|_{L^{q'}(\mathbb{R}^{n-1})}\tag{4.6.6}
\label{eq:4.6.6}\end{equation}
当 $p' = \frac{n+1}{n-1}\frac{q'}{q'-1}$ 时,在
(1) $1 \leq q' \leq 2$,$n \geq 3$,
(2) $1 \leq q' < 4$,$n = 2$
的情形下成立.
\end{proposition}

\begin{proof}[命题~\ref{prop:4.6.2} (1) 的证明]
$(q', p') = (1, \infty)$ 的情形显然,故证明 $(q', p') = \left(2, \frac{2(n+1)}{n-1}\right)$ 的情形:
\begin{equation}
\|T_{\lambda}g\|_{L^{\frac{2(n+1)}{n-1}}(\mathbb{R}^{n})} \leq C\lambda^{-\frac{n(n-1)}{2(n+1)}}\|g\|_{L^{2}(\mathbb{R}^{n-1})},\tag{4.6.7}
\label{eq:4.6.7}\end{equation}
再用插值法导出结论.

考虑对偶性,为证明 \eqref{eq:4.6.7},令
\begin{equation*}
T_{\lambda}^{*}f(y) := \int_{\mathbb{R}^{n}}e^{i\lambda\phi(z, y)}a(z, y)f(z)\dd z,
\end{equation*}
证明
\begin{equation}
\|T_{\lambda}^{*}f\|_{L^{2}(\mathbb{R}^{n-1})} \leq C\lambda^{-\frac{n(n-1)}{2(n+1)}}\|f\|_{L^{\frac{2(n+1)}{n+3}}(\mathbb{R}^{n})}.\tag{4.6.8}
\label{eq:4.6.8}\end{equation}
此左端可以直接写出:
\begin{equation}
\begin{aligned}
\|T_{\lambda}^{*}f\|_{L^{2}(\mathbb{R}^{n-1})}^{2} & = \int_{\mathbb{R}^{n-1}}\int_{\mathbb{R}^{n}}\int_{\mathbb{R}^{n}}e^{i\lambda[\phi(z, y) - \phi(w, y)]}a(z, y)\overline{a(w, y)}f(z)\overline{f(w)}\dd z\dd w\dd y \\
& \leq \|\tilde{K}_{\lambda}f\|_{L^{\frac{2(n+1)}{n-1}}(\mathbb{R}^{n})}\times\|f\|_{L^{\frac{2(n+1)}{n+3}}(\mathbb{R}^{n})}.
\end{aligned}
\end{equation}
这里
\begin{equation}
\begin{aligned}
\tilde{K}_{\lambda}f(z) & := T_{\lambda}\overline{T_{\lambda}^{*}f}(z) = \int_{\mathbb{R}^{n}}K_{\lambda}(z, w)\overline{f(w)}\dd w, \\
K_{\lambda}(z, w) & := \int_{\mathbb{R}^{n-1}}e^{i\lambda[\phi(z, y) - \phi(w, y)]}a(z, y)\overline{a(w, y)}\dd y.
\end{aligned}
\end{equation}
因此,为证明 \eqref{eq:4.6.8},令
\begin{equation*}
K_{\lambda}f(z) := \int_{\mathbb{R}^{n}}K_{\lambda}(z, w)f(w)\dd w,
\end{equation*}
当 $r = \frac{2(n+1)}{n+3}$ 时证明
\begin{equation}
\|K_{\lambda}f\|_{L^{r'}(\mathbb{R}^{n})} \leq C\lambda^{-\frac{n(n-1)}{n+1}}\|f\|_{L^{r}(\mathbb{R}^{n})}.\tag{4.6.9}
\label{eq:4.6.9}\end{equation}
这个式子的有用之处在于,两边都能估计同一维数的函数.

将 $z = (z', u)$($z' \in \mathbb{R}^{n-1}$),$w = (w', v)$($w' \in \mathbb{R}^{n-1}$)分别分割,固定 $u, v$,令
\begin{equation}
\begin{aligned}
K_{u,v}^{\lambda}f(z') & := \int_{\mathbb{R}^{n-1}}K_{u,v}^{\lambda}(z', w')f(w', v)\dd w', \\
K_{u,v}^{\lambda}(z', w') & := K_{\lambda}(z', u; w', v),
\end{aligned}
\end{equation}
若证明
\begin{equation}
\|K_{u,v}^{\lambda}f\|_{L^{r'}(\mathbb{R}^{n-1})} \leq C\lambda^{-\frac{n(n-1)}{n+1}}|u - v|^{-1+(1/r - 1/r')}\|f(\cdot, v)\|_{L^{r}(\mathbb{R}^{n-1})},\tag{4.6.10}
\label{eq:4.6.10}\end{equation}
则经由一维 Hardy--Littlewood--Sobolev 不等式到达 \eqref{eq:4.6.9}.

实际上,由于
\begin{equation*}
K_{u}^{\lambda}f(z') := \int_{\mathbb{R}^{n}}K_{\lambda}(z', u; w)f(w)\dd w = \int_{-\infty}^{\infty}K_{u,v}^{\lambda}f(z')\dd v,
\end{equation*}
由 Minkowski 积分不等式,
\begin{equation}
\begin{aligned}
\|K_{u}^{\lambda}f\|_{L^{r'}(\mathbb{R}^{n-1})} & \leq \int_{-\infty}^{\infty}\|K_{u,v}^{\lambda}f\|_{L^{r'}(\mathbb{R}^{n-1})}\dd v \\
& \leq C\lambda^{-\frac{n(n-1)}{n+1}}\int_{-\infty}^{\infty}|u - v|^{-1+(1/r - 1/r')}\|f(\cdot, v)\|_{L^{r}(\mathbb{R}^{n-1})}\dd v.
\end{aligned}
\end{equation}
上式左端关于 $u$ 的 $L^{r'}(\mathbb{R})$ 范数等于 \eqref{eq:4.6.9} 的左端,故应用 Hardy--Littlewood--Sobolev 不等式(定理~\ref{thm:4.2.3}),
\begin{equation}
\begin{aligned}
\|K_{\lambda}f\|_{L^{r'}(\mathbb{R}^{n})} & \leq C\lambda^{-\frac{n(n-1)}{n+1}}\left\|\int_{-\infty}^{\infty}|u - v|^{-1+(1/r - 1/r')}\|f(\cdot, v)\|_{L^{r}(\mathbb{R}^{n-1})}\dd v\right\|_{L^{r'}(\mathbb{R})} \\
& \leq C\lambda^{-\frac{n(n-1)}{n+1}}\left\|\|f(\cdot, v)\|_{L^{r}(\mathbb{R}^{n-1})}\right\|_{L^{r}(\mathbb{R})} \\
& = C\lambda^{-\frac{n(n-1)}{n+1}}\|f\|_{L^{r}(\mathbb{R}^{n})}.
\end{aligned}
\end{equation}

\eqref{eq:4.6.10} 由 $L^{2} \to L^{2}$,$L^{1} \to L^{\infty}$ 的插值法得到.

前者,将
\begin{equation}
K_{u,v}^{\lambda}f(z') = \int_{\mathbb{R}^{n-1}}e^{i\lambda\phi(z, y)}a(z, y)\int_{\mathbb{R}^{n-1}}\overline{e^{i\lambda\phi(w, y)}a(w, y)}f(w)\dd w'\dd y
\end{equation}
还原,并想起
\begin{equation*}
\phi(w, y) = w_{1}y_{1} + \ldots + w_{n-1}y_{n-1} + uh(y),
\end{equation*}
则
\begin{equation*}
\det\left(\left(\partial^{2}/\partial w_{j}\partial y_{k}\right)\phi(w, y)\right)_{j, k \leq n-1} = 1 \neq 0,
\end{equation*}
故应用定理~\ref{thm:4.5.1},
\begin{equation}
\begin{aligned}
& \left\|\int_{\mathbb{R}^{n-1}}\overline{e^{i\lambda\phi(w, y)}a(w, y)}f(w)\dd w'\right\|_{L^{2}(\mathbb{R}^{n-1})} \\
& \leq C\lambda^{-(n-1)/2}\|f(\cdot, v)\|_{L^{2}(\mathbb{R}^{n-1})}.
\end{aligned}
\end{equation}
同样考虑,
\begin{equation}
\|K_{u,v}^{\lambda}f\|_{L^{2}(\mathbb{R}^{n-1})} \leq C\lambda^{-(n-1)}\|f(\cdot, v)\|_{L^{2}(\mathbb{R}^{n-1})}.\tag{4.6.11}
\label{eq:4.6.11}\end{equation}

用驻相法证明后者.需要估计的核
\begin{equation}
K_{u,v}^{\lambda}(z', w') = \int_{\mathbb{R}^{n-1}}e^{i\lambda\Phi(y, z, w)}a(z, y)\overline{a(w, y)}\dd y,
\end{equation}
\begin{equation*}
\Phi(y, z, w) := (z_{1} - w_{1})y_{1} + \ldots + (z_{n-1} - w_{n-1})y_{n-1} + (u - v)h(y),
\end{equation*}
验证
\begin{equation}
|K_{u,v}^{\lambda}(z', w')| \leq C(\lambda|u - v|)^{-(n-1)/2}.\tag{4.6.12}
\label{eq:4.6.12}\end{equation}

将
\begin{equation}
\begin{aligned}
& (z_{1} - w_{1})y_{1} + \ldots + (z_{n-1} - w_{n-1})y_{n-1} + (u - v)h(y) \\
& = (u - v)\left(\frac{z_{1} - w_{1}}{u - v}y_{1} + \ldots + \frac{z_{n-1} - w_{n-1}}{u - v}y_{n-1} + h(y)\right)
\end{aligned}
\end{equation}
中的 $\mathbb{R}^{n}$ 上的点设为
\begin{equation*}
A(z, w) := \frac{1}{u - v}(z_{1} - w_{1}, \ldots, z_{n-1} - w_{n-1}, 1),
\end{equation*}
则可知,仅当连接原点与 $A(z, w)$ 的直线与 $\tilde{D}$ 有交点 $(y_{0}, h(y_{0}))$ 时,
\begin{equation*}
\nabla_{y}\Phi(y_{0}, w, z) = 0.
\end{equation*}
不失一般性,可设
\begin{equation*}
(y_{0}, h(y_{0})) = A(z_{0}, w_{0}) = (0, \ldots, 0, 1)
\end{equation*}
继续验证.此时
\begin{equation*}
\Phi(y, z_{0}, w_{0}) = h(y),
\end{equation*}
故
\begin{equation*}
\det\left(\left(\partial^{2}/\partial y_{j}\partial y_{k}\right)\Phi(y, z_{0}, w_{0})\right) = 1 \neq 0 \quad (y = 0),
\end{equation*}
由定理~\ref{thm:4.4.2} 可估计为
\begin{equation*}
\left|K_{u,v}^{\lambda}(0, z_{0}, w_{0})\right| \leq C(\lambda|u - v|)^{-(n-1)/2}.
\end{equation*}
其他情形得到比这更好的估计,作为结论 \eqref{eq:4.6.12} 得证.由此得到
\begin{equation}
\|K_{u,v}^{\lambda}f\|_{L^{\infty}(\mathbb{R}^{n-1})} \leq C\lambda^{-(n-1)/2}|u - v|^{-(n-1)/2}\|f(\cdot, v)\|_{L^{1}(\mathbb{R}^{n-1})}.\tag{4.6.13}
\label{eq:4.6.13}\end{equation}

利用
\begin{equation*}
\begin{cases}
\frac{1}{r} = \frac{1}{2}\left(1 - \frac{2}{n+1}\right) + \frac{2}{n+1}, \\
-\frac{n(n-1)}{n+1} = -(n-1)\left(1 - \frac{2}{n+1}\right) - \frac{n-1}{2}\frac{2}{n+1}, \\
-1 + \left(1/r - 1/r'\right) = -\frac{n-1}{n+1} = -\frac{n-1}{2}\frac{2}{n+1},
\end{cases}
\end{equation*}
插值 \eqref{eq:4.6.11},\eqref{eq:4.6.13},\eqref{eq:4.6.10} 得证.
\end{proof}

\begin{proof}[命题~\ref{prop:4.6.2} (2) 的证明]
这里稍微改变形式,令
\begin{equation}
\begin{aligned}
T_{\lambda}g(z) & := \int_{-\infty}^{\infty}e^{i\lambda\phi(z, \theta)}a(z, \theta)g(\theta)\dd\theta, \\
\phi(z, \theta) & := z_{1}\cos\theta + z_{2}\sin\theta,
\end{aligned}
\end{equation}
在 $p' = 3\frac{q'}{q'-1}$,$1 \leq q' < 4$ 中证明
\begin{equation}
\|T_{\lambda}g\|_{L^{p'}(\mathbb{R}^{2})} \leq C\lambda^{-2/p'}\|g\|_{L^{q'}(\mathbb{R})}.\tag{4.6.14}
\label{eq:4.6.14}\end{equation}
证明的关键在于 $p' > 4$.在这种情况下,可以估计
\begin{equation}
(T_{\lambda}g(z))^{2} = \int_{\mathbb{R}^{2}}e^{i\lambda\Phi(z, y)}b(z, y)g(y_{1})g(y_{2})\dd y,
\end{equation}
\begin{equation*}
\Phi(z, y) := z_{1}(\cos y_{1} + \cos y_{2}) + z_{2}(\sin y_{1} + \sin y_{2}), \quad b(z, y) := a(z, y_{1})a(z, y_{2}).
\end{equation*}

想对
\begin{equation*}
\int_{\mathbb{R}^{2}}e^{i\lambda\Phi(z, y)}b(z, y)f(y)\dd y
\end{equation*}
应用定理~\ref{thm:4.5.1}.但是,由于
\begin{equation*}
\left|\begin{array}{cc}
(\partial^{2}/\partial z_{1}\partial y_{1})\Phi(z, y) & (\partial^{2}/\partial z_{1}\partial y_{2})\Phi(z, y) \\
(\partial^{2}/\partial z_{2}\partial y_{1})\Phi(z, y) & (\partial^{2}/\partial z_{2}\partial y_{2})\Phi(z, y)
\end{array}\right| = -\sin(y_{1} - y_{2}),
\end{equation*}
不能这样直接应用.于是,变换坐标
\begin{equation*}
(u_{1}, u_{2}) = (y_{1} - y_{2}, y_{1} + y_{2}),
\end{equation*}
则 $|du/dy| = 2$,故必要时将 $b$ 的支集分解小,可得
\begin{equation*}
|\det(\partial^{2}\Phi/\partial z\partial u)| \geq c|u_{1}|.
\end{equation*}
进而
\begin{equation}
\begin{aligned}
\Phi(z, u) & = z_{1}\left(\cos\frac{u_{2} + u_{1}}{2} + \cos\frac{u_{2} - u_{1}}{2}\right) \\
& \quad + z_{2}\left(\sin\frac{u_{2} + u_{1}}{2} + \sin\frac{u_{2} - u_{1}}{2}\right),
\end{aligned}
\end{equation}
\begin{equation*}
b(z, u) = a\left(z, \frac{u_{2} + u_{1}}{2}\right)a\left(z, \frac{u_{2} - u_{1}}{2}\right),
\end{equation*}
故 $\Phi(z, u)$,$b(z, u)$ 都是 $u_{1}$ 的偶函数,可以看作 $u_{1}^{2}$ 的连续函数.于是,变换坐标
\begin{equation*}
(v_{1}, v_{2}) = \left(\frac{1}{2}u_{1}^{2}, u_{2}\right),
\end{equation*}
则 $|dv/du| = |u_{1}|$,故
\begin{equation*}
|\det(\partial^{2}\Phi/\partial z\partial v)| \geq c
\end{equation*}
成立.这里,对
\begin{equation*}
\tilde{T}_{\lambda}f(z) := \int_{\mathbb{R}^{2}}e^{i\lambda\Phi(z, v)}b(z, v)f(v)\dd v
\end{equation*}
应用定理~\ref{thm:4.5.1},特别地系~\ref{cor:4.5.2} 成立.即
\begin{equation}
\|\tilde{T}_{\lambda}f\|_{L^{r'}(\mathbb{R}^{2})} \leq C\lambda^{-2/r'}\|f\|_{L^{r}(\mathbb{R}^{2})} \quad (1 \leq r \leq 2).\tag{4.6.15}
\label{eq:4.6.15}\end{equation}

由定义可计算为
\begin{equation}
\begin{aligned}
\|T_{\lambda}g\|_{L^{p'}(\mathbb{R}^{2})}^{2} & = \|(T_{\lambda}g)^{2}\|_{L^{p'/2}(\mathbb{R}^{2})} \\
& = \|\tilde{T}_{\lambda}g(\cdot)g(\cdot)|dv/dy|^{-1}\|_{L^{p'/2}(\mathbb{R}^{2})},
\end{aligned}
\end{equation}
对右端应用 \eqref{eq:4.6.15},
\begin{equation}
\begin{aligned}
\|T_{\lambda}g\|_{L^{p'}(\mathbb{R}^{2})}^{2} & \leq C\lambda^{-4/p'}\left(\int_{\mathbb{R}^{2}}g(y_{1})^{r}g(y_{2})^{r}|y_{1} - y_{2}|^{1-r}\dd y\right)^{1/r}, \\
& \qquad r = p'/(p' - 2).
\end{aligned}
\end{equation}
这里,可计算出
\begin{equation*}
1 - r = -1 + (2 - r), \quad 2 - r = [q'/r]^{-1} - [(q'/r)']^{-1},
\end{equation*}
用一维 Hardy--Littlewood--Sobolev 不等式到达结论.
\end{proof}
